% ============================================================
% Paper deep reading: E-GRPO
% Source: https://arxiv.org/html/2601.00423v1
% ============================================================

\section{论文精读：E-GRPO}

\begin{conceptbox}
\textbf{论文：} \href{https://arxiv.org/abs/2601.00423}{E-GRPO: High Entropy Steps Drive Effective Reinforcement Learning for Flow Models}\\
\textbf{作者：} Shengjun Zhang, Zhang Zhang, Chensheng Dai, Yueqi Duan。\\
\textbf{简称：} E-GRPO。\\
\textbf{版本：} arXiv 2601.00423v1。\\
\textbf{代码：} \href{https://github.com/shengjun-zhang/VisualGRPO}{VisualGRPO}。\\
\textbf{方向标签：} Flow Matching、GRPO、Entropy-aware Exploration、Step Merging、Visual Preference Alignment。\\
\textbf{阅读日期：} 2026-05-06。\\
\textbf{阅读口径：} 这篇和 OP-GRPO 都属于 ``sample efficiency / exploration efficiency'' 支线，但切入点不同：OP-GRPO 问 \textbf{旧轨迹能不能复用}，E-GRPO 问 \textbf{哪些 denoising steps 真正值得随机探索和做 RL 更新}。
\end{conceptbox}

\subsection{0. 一句话结论}

E-GRPO 的核心判断很直接：Flow-GRPO 不应该在所有 denoising steps 上均匀做 SDE 探索和 GRPO 更新，因为不同时间步的探索熵不一样。

\[
\begin{aligned}
\text{E-GRPO}
&= \text{entropy-aware step selection}\\
&\quad + \text{low-entropy step merging}\\
&\quad + \text{multi-step group-normalized advantage}
\end{aligned}
\]

它认为 high-entropy steps 会产生更有区分度的 rollouts，更适合用 reward 做相对比较；low-entropy steps 的随机扰动太小，生成结果差异不明显，继续把这些 steps 当成独立 RL action 只会增加计算和 reward attribution 噪声。

\begin{summarybox}
\textbf{我的判断：} E-GRPO 不是在改 reward，也不是在做 replay buffer，而是在改 \key{探索发生的位置}。它把 Flow-GRPO 里的 ``每一步都随机、每一步都更新'' 改成 ``只在高信息量的 consolidated SDE step 上随机和更新''。这使它和 OP-GRPO 正好互补：一个提高轨迹复用效率，一个提高单次 rollout 内部的步骤效率。
\end{summarybox}

\subsection{1. 它为什么重要}

Flow-GRPO 把 deterministic flow ODE 改写成 stochastic SDE policy，从而能计算 transition probability、做 GRPO-style policy update。这个改写解决了 ``flow model 没有 logprob 怎么做 RL'' 的问题，但也引入了一个新的默认假设：

\[
\text{所有 denoising steps 都值得注入随机性，也都值得接受同样形式的 RL 更新。}
\]

E-GRPO 认为这个假设是有问题的。因为在 flow / diffusion 采样里，不同时间步的随机扰动影响完全不同：

\begin{itemize}[leftmargin=1.5em]
  \item 早期 high-noise / high-entropy steps 会改变整体构图、语义布局和主体形态，生成结果之间差异大。
  \item 后期 low-noise / low-entropy steps 更多是在局部细节附近微调，rollouts 之间区分度小。
\end{itemize}

如果后期低熵步骤产生的几条样本几乎看不出差异，但我们还强行对它们做 group-relative reward normalization，那么 reward 里的随机波动、scorer 噪声、下游步骤误差就会被误当成该 step 的学习信号。

\begin{warnbox}
\textbf{关键问题：} Flow-GRPO 的稀疏 reward 不只是不知道 ``哪一步贡献大''，还不知道 ``哪一步本来就没有足够探索空间''。如果某一步的 rollout 差异太小，组内相对优势本身就很难可靠。
\end{warnbox}

所以 E-GRPO 的问题可以概括成：

\begin{quote}
在 flow-model RL 中，怎样根据 step entropy 选择真正有探索价值的时间步，并避免低熵多步随机性带来的 reward attribution ambiguity？
\end{quote}

\subsection{2. 熵视角：为什么不是所有 SDE steps 都一样}

论文从 Flow-GRPO 的 SDE 采样形式出发。为了让 flow model 可以探索，原本的 ODE update 会被改写成带随机性的 SDE transition：

\[
x_{t+\Delta t}
=
x_t
+
\mu_{\theta}(x_t,t)\Delta t
+
\sigma_t\sqrt{\Delta t}\epsilon,
\qquad
\epsilon\sim\mathcal{N}(0,I).
\]

其中论文使用的噪声系数形式是：

\[
\sigma_t
=
a\sqrt{\frac{t}{1-t}}.
\]

对于反向 SDE step，它的 transition distribution 可以看成 Gaussian，因此熵主要由 covariance 控制。论文给出的 step entropy 可写成：

\[
h(t)
=
\frac{d}{2}\log\left(2\pi e\sigma_t^2\Delta t\right).
\]

代入 \(\sigma_t=a\sqrt{t/(1-t)}\) 后，有：

\[
h(t)
=
\frac{d}{2}
\log\left(
2\pi e\cdot a^2\cdot \frac{t}{1-t}\cdot \Delta t
\right).
\]

也就是说，在忽略常数后：

\[
e^{h(t)}
\propto
\frac{t}{1-t}\Delta t.
\]

这给了 E-GRPO 一个很明确的判断标准：哪个 step 的 SDE entropy 更大，哪个 step 的随机探索范围就更大，也更可能产生 reward 上可区分的 rollouts。

\begin{intubox}
\textbf{直觉版：} high-entropy step 像是在问 ``我要往哪个语义方向走''；low-entropy step 更像是在问 ``边缘再抹平一点还是锐一点''。前者更适合 RL 探索，后者如果也用同样的 GRPO 强度去训，收益小、噪声大。
\end{intubox}

\begin{figure}[htbp]
\centering
\includegraphics[width=0.95\linewidth]{egrpo/01_entropy_influence.png}
\caption{E-GRPO 对 step entropy 的经验观察。左侧比较不同 SDE sampling 策略下的生成差异，中间展示合并步数如何提升 entropy，右侧展示只训练 high-entropy early steps 的 reward 曲线更优。图源：\href{https://arxiv.org/html/2601.00423v1}{E-GRPO arXiv HTML Figure 1}。}
\label{fig:egrpo-entropy-influence}
\end{figure}

\subsection{3. 方法一：Entropy-driven Step Merging}

\subsubsection{3.1 用阈值区分 high-entropy 和 low-entropy steps}

E-GRPO 首先为每个 denoising timestep 计算 entropy proxy，然后用阈值 \(\tau\) 将时间步分成两类：

\[
\begin{aligned}
\text{high-entropy steps} &: e^{h(t_k)}\geq \tau,\\
\text{low-entropy steps} &: e^{h(t_k)}<\tau.
\end{aligned}
\]

实际 flow sampling 里通常还会做 time shift：

\[
t_k
=
\frac{s\hat{t}_k}{1+(s-1)\hat{t}_k},
\qquad
\hat{t}_k=\frac{k}{T}.
\]

论文将这个时间变换代入 entropy proxy 后，用它来判断哪些 steps 具有足够探索熵。

\subsubsection{3.2 低熵步骤不单独做 SDE，而是合并成一个更大的 SDE step}

对连续低熵步骤 \(\{t_m,\ldots,t_{m-l}\}\)，E-GRPO 不再逐步注入 SDE noise，而是把它们合并成一个 consolidated SDE step：

\[
x_{t_{m-l}}
=
x_{t_m}
+
\mu_{\theta}(x_{t_m},t_m)(t_m-t_{m-l})
+
\sigma_t\sqrt{t_m-t_{m-l}}\epsilon.
\]

合并后的 entropy 近似满足：

\[
e^{h(t_k)}
\propto
\frac{t_m}{1-t_m}(t_m-t_{m-l}).
\]

因为 \((t_m-t_{m-l})\) 变大，所以多个低熵小 step 可以合成一个达到阈值附近的高熵 step。

\begin{summarybox}
\textbf{这一步的设计很关键：} E-GRPO 不是简单跳过 low-entropy steps，而是把连续 low-entropy steps 合并成一个 ``刚好有足够探索空间'' 的 SDE event。这样既减少了独立随机步骤数量，又避免完全失去这段轨迹上的探索能力。
\end{summarybox}

\begin{figure}[htbp]
\centering
\includegraphics[width=0.82\linewidth]{egrpo/02_pipeline.png}
\caption{E-GRPO 的整体 sampling strategy：先为不同 timesteps 构造 anchor noise latents，再根据 entropy analysis 为每个 active SDE timestep 选择 merged steps，最后在对应 group 内生成结果并计算 advantage。图源：\href{https://arxiv.org/html/2601.00423v1}{E-GRPO arXiv HTML Figure 2}。}
\label{fig:egrpo-pipeline}
\end{figure}

\subsubsection{3.3 为什么不用固定 2-step / 4-step / 6-step merge}

一个简单做法是固定每 \(l\) 步合并一次。但 E-GRPO 的消融显示，固定合并不如 adaptive merging。

原因也很自然：不同 timestep 的 entropy 本来就不一样。如果固定 4-step merge，有些位置合并后还是太低熵，有些位置可能熵过高，导致探索方向太散。论文因此为每个 low-entropy timestep 自适应选择 \(l\)，让合并后的 entropy 刚好超过阈值 \(\tau\)。

\[
\text{choose }l
\quad\text{s.t.}\quad
 e^{h(t_m, l)} \gtrsim \tau.
\]

这相当于把每个 consolidated SDE step 的探索尺度校准到一个相近范围。

\subsection{4. 方法二：Multi-step Group Normalized Advantage}

合并步骤之后，advantage 也不能再按原来的单步方式理解。

对一个 active SDE timestep \(t_n\)，它对应一个合并后的时间段：

\[
\mathcal{T}_n
=
\{t_n,t_{n-1},\ldots,t_{n-l_n}\}.
\]

给定同一个 prompt \(c\)，论文为每个 active timestep \(t_n\) 生成 \(G^{(n)}\) 条 trajectories，并只在共享同一个 consolidated SDE step 的样本内部计算 group-normalized advantage：

\[
A_{t_n}^{(i)}
=
\frac{
R(x_{0,t_n}^{(i)},c)
-
\mathrm{mean}\{R(x_{0,t_n}^{(j)},c)\}_{j=1}^{G^{(n)}}
}{
\mathrm{std}\{R(x_{0,t_n}^{(j)},c)\}_{j=1}^{G^{(n)}}
}.
\]

对应的 ratio 是合并 step 的 transition ratio：

\[
r_{t_n}^{(i)}(\theta)
=
\frac{
p_{\theta}(x_{t_{n-l_n}}^{(i)}\mid x_{t_n}^{(i)},c)
}{
p_{\theta_{\mathrm{old}}}(x_{t_{n-l_n}}^{(i)}\mid x_{t_n}^{(i)},c)
}.
\]

于是 clipped surrogate 仍然保持 GRPO / PPO 风格：

\[
\frac{1}{\hat{T}}
\sum_{n=N}^{T}
\frac{1}{G^{(n)}}
\sum_{i=1}^{G^{(n)}}
\min\left(
 r_{t_n}^{(i)}A_{t_n}^{(i)},
 \mathrm{clip}(r_{t_n}^{(i)},1-\epsilon,1+\epsilon)A_{t_n}^{(i)}
\right).
\]

\begin{figure}[htbp]
\centering
\includegraphics[width=0.72\linewidth]{egrpo/03_ambiguous_reward.png}
\caption{连续多步 SDE sampling 会让同一个 final advantage 同时对应多个 timesteps，从而产生 ambiguous reward signal；E-GRPO 通过 merged-step SDE sampling 把最终 advantage 对齐到一个 merged SDE step，降低单步 credit assignment 的歧义。图源：\href{https://arxiv.org/html/2601.00423v1}{E-GRPO arXiv HTML Figure 3}。}
\label{fig:egrpo-ambiguous-reward}
\end{figure}

\begin{warnbox}
\textbf{这里和 DenseGRPO 不一样：} DenseGRPO 是想给每个 step 估计更密的 reward gain；E-GRPO 是先把低信息量 steps 合成一个有意义的 exploration event，再在这个 event 内做 group-relative advantage。它不是 ``更细''，而是 ``该细的地方细，该合的地方合''。
\end{warnbox}

\subsection{5. 算法流程：训练时到底怎么做}

可以把 E-GRPO 的训练流程压缩成五步：

\begin{enumerate}[leftmargin=1.5em]
  \item 计算所有 timesteps 的 entropy proxy \(h(t_k)\)。
  \item 根据阈值 \(\tau\) 识别 high-entropy steps 和 low-entropy steps。
  \item 对连续 low-entropy steps 自适应选择 \(l_n\)，形成 consolidated SDE steps。
  \item 对非 active steps 使用 ODE sampling，对 consolidated active steps 使用 SDE sampling。
  \item 在每个 merged-step group 内计算 advantage 和 clipped GRPO objective。
\end{enumerate}

因此，E-GRPO 实际上把 rollout policy 改成了一个 ODE/SDE hybrid：

\[
\text{ODE on uninformative steps}
\quad + \quad
\text{SDE on entropy-calibrated merged steps}.
\]

\begin{intubox}
\textbf{直觉：} 它不是减少采样步数本身，而是减少 ``无效随机性'' 和 ``无效 RL 更新''。ODE steps 负责稳定推进生成，SDE merged steps 负责真正探索。
\end{intubox}

\subsection{6. 实验设置}

论文的实验设置如下：

\begin{itemize}[leftmargin=1.5em]
  \item \textbf{Backbone：} FLUX.1-dev。
  \item \textbf{数据集：} HPD，约 103K training prompts 和 400 testing prompts。
  \item \textbf{采样步数：} 总采样步数 \(T=16\)。
  \item \textbf{噪声参数：} \(a=0.7\)。
  \item \textbf{默认 entropy threshold：} \(\tau=2.2\)。
  \item \textbf{训练步数：} 300 iterations。
  \item \textbf{训练 reward：} 先使用 HPS-v2.1；为了缓解 reward hacking，又使用 HPS-v2.1 + CLIP 的联合 reward。
  \item \textbf{评测指标：} HPS、CLIP、PickScore、ImageReward。
\end{itemize}

这套设置有一个值得注意的点：作者明确指出只用 HPS 训练可能导致 reward hacking，所以主结果里同时报告了单 reward 和 HPS+CLIP 两种训练方式。

\subsection{7. 主要结果：E-GRPO 的收益在哪里}

\subsubsection{7.1 单 HPS reward 下，HPS 指标最好}

在只用 HPS-v2.1 作为训练 reward 时，E-GRPO 在 HPS 上达到 \textbf{0.391}，高于几个主要 baseline：

\begin{center}
\resizebox{\linewidth}{!}{
\begin{tabular}{lcccc}
\toprule
\textbf{方法} & \textbf{HPS} & \textbf{CLIP} & \textbf{PickScore} & \textbf{ImageReward}\\
\midrule
FLUX.1-dev & 0.311 & 0.388 & 0.231 & 1.089\\
DanceGRPO & 0.353 & 0.375 & 0.228 & 1.233\\
MixGRPO & 0.378 & 0.358 & 0.225 & 1.266\\
GranularGRPO & 0.385 & 0.355 & 0.229 & 1.313\\
TempFlowGRPO & 0.382 & 0.357 & 0.231 & 1.264\\
E-GRPO & \textbf{0.391} & 0.355 & \textbf{0.232} & \textbf{1.324}\\
\bottomrule
\end{tabular}
}
\end{center}

论文强调，E-GRPO 相比 DanceGRPO 在 HPS 上提升约 \textbf{10.8\%}。但我更关注另一个现象：HPS 提高的同时，CLIP 并没有同步提高，说明单 reward 训练确实存在向 HPS scorer 偏移的风险。

\subsubsection{7.2 HPS+CLIP 联合 reward 下，泛化指标更稳}

在 HPS+CLIP 联合训练下，E-GRPO 的表现更有说服力：

\begin{center}
\resizebox{\linewidth}{!}{
\begin{tabular}{lcccc}
\toprule
\textbf{方法} & \textbf{HPS} & \textbf{CLIP} & \textbf{PickScore} & \textbf{ImageReward}\\
\midrule
FLUX.1-dev & 0.311 & 0.388 & 0.231 & 1.089\\
DanceGRPO & 0.331 & 0.389 & 0.227 & 1.128\\
MixGRPO & 0.363 & 0.399 & 0.230 & 1.436\\
GranularGRPO & 0.377 & 0.400 & 0.236 & 1.490\\
BranchGRPO & 0.342 & 0.384 & 0.230 & 1.243\\
TempFlowGRPO & 0.310 & 0.388 & 0.230 & 1.106\\
E-GRPO & \textbf{0.382} & \textbf{0.401} & \textbf{0.237} & \textbf{1.494}\\
\bottomrule
\end{tabular}
}
\end{center}

论文报告，相比 DanceGRPO，E-GRPO 在 ImageReward 上提升 \textbf{32.4\%}，在 PickScore 上提升 \textbf{4.4\%}。这说明 entropy-guided step selection 不只是刷某一个训练 reward，而是能带来更广泛的人类偏好指标改善。

\subsubsection{7.3 训练曲线：早期增长更快，收敛更平滑}

论文还展示了 reward trajectory：E-GRPO 早期 reward 增长更快，最终 reward 更高，曲线也更平滑。

这和方法假设是一致的：如果训练集中在 high-entropy / consolidated informative steps 上，那么每次更新都更可能改变真正影响 reward 的生成方向，而不是在低熵后期步骤上浪费梯度。

\begin{figure}[htbp]
\centering
\includegraphics[width=0.95\linewidth]{egrpo/04_reward_curve.png}
\caption{E-GRPO 与 baseline 的训练 reward 曲线对比。曲线说明 entropy-guided exploration 能更快提升训练 reward，并且相比全 timestep 更新更稳定。图源：\href{https://arxiv.org/html/2601.00423v1}{E-GRPO arXiv HTML Figure 4}。}
\label{fig:egrpo-reward-curve}
\end{figure}

\subsubsection{7.4 可视化结果：语义和细节更均衡}

论文的可视化对比显示，E-GRPO 相比 baseline 更容易同时保留 prompt 语义和细粒度视觉细节。这个现象和前面的机制解释是一致的：high-entropy steps 负责更大的语义探索，merged-step SDE 又避免后期低熵步骤把 reward 噪声放大成错误更新。

\begin{figure}[htbp]
\centering
\includegraphics[width=0.95\linewidth]{egrpo/05_visual_comparison.png}
\caption{E-GRPO 与其他 baseline 的生成结果可视化对比。论文认为 E-GRPO 能更好地整合语义对齐和细节质量。图源：\href{https://arxiv.org/html/2601.00423v1}{E-GRPO arXiv HTML Figure 5}。}
\label{fig:egrpo-visual-comparison}
\end{figure}

\subsection{8. 消融：哪些设计真的重要}

\subsubsection{8.1 Adaptive merging 优于固定步数 merging}

论文比较了固定 2-step、4-step、6-step merging 和 adaptive merging：

\begin{center}
\resizebox{\linewidth}{!}{
\begin{tabular}{lcccc}
\toprule
\textbf{Merging strategy} & \textbf{HPS} & \textbf{CLIP} & \textbf{PickScore} & \textbf{ImageReward}\\
\midrule
2-step & 0.382 & 0.290 & 0.232 & 1.223\\
4-step & 0.374 & 0.302 & 0.230 & 1.216\\
6-step & 0.337 & \textbf{0.372} & 0.226 & 1.298\\
Adaptive & \textbf{0.391} & 0.355 & \textbf{0.232} & \textbf{1.324}\\
\bottomrule
\end{tabular}
}
\end{center}

这说明固定 merge 的问题在于不看局部 entropy：合并太少，探索仍然不足；合并太多，单个 SDE event 过于随机，reward 方向又会变模糊。Adaptive merging 的优势就是把每个 active step 的探索强度拉回到一个合适范围。

\subsubsection{8.2 只训前 8 个 high-entropy steps 最好}

论文还比较了训练不同 timestep subset：

\begin{center}
\resizebox{\linewidth}{!}{
\begin{tabular}{lcccc}
\toprule
\textbf{训练 steps} & \textbf{HPS} & \textbf{CLIP} & \textbf{PickScore} & \textbf{ImageReward}\\
\midrule
First 4 Steps & 0.370 & 0.348 & 0.231 & 1.252\\
First 8 Steps & \textbf{0.391} & 0.355 & \textbf{0.232} & \textbf{1.324}\\
Second 8 Steps & 0.357 & \textbf{0.381} & 0.231 & 1.250\\
Full Steps & 0.366 & 0.359 & 0.231 & 1.169\\
\bottomrule
\end{tabular}
}
\end{center}

结论非常清楚：前 8 个 high-entropy steps 已经足以获得最强性能；训练所有 steps 不但成本更高，效果还不如只训前半段；只训后 8 个 low-entropy steps 则明显更弱。

\begin{summarybox}
\textbf{我会这样记这组消融：} E-GRPO 不是 ``少训一些 steps 也能省钱'' 这么简单，而是证明了 \key{后半段低熵 steps 的 RL 信号质量确实差}。全步骤训练把有用信号和低质量信号混在一起，反而拖累优化。
\end{summarybox}

\subsection{9. 和 OP-GRPO / DenseGRPO / TempFlow-GRPO 的关系}

把 E-GRPO 放进这条 Flow-GRPO 后训练主线里，它的位置大概是：

\begin{center}
\resizebox{\linewidth}{!}{
\begin{tabular}{lll}
\toprule
\textbf{论文} & \textbf{主要改哪里} & \textbf{核心问题}\\
\midrule
TempFlow-GRPO & temporal weighting / branching & 哪个 timestep 更该分到 credit\\
DenseGRPO & dense reward gain & terminal reward 怎样分给每个 step\\
OP-GRPO & off-policy replay & 已经采出来的高质量 trajectory 能否复用\\
E-GRPO & entropy-aware exploration & 哪些 steps 值得随机探索和 RL 更新\\
\bottomrule
\end{tabular}
}
\end{center}

\begin{itemize}[leftmargin=1.5em]
  \item \textbf{和 OP-GRPO：} 两者都关心效率，但 OP-GRPO 是跨 iteration 的样本复用，E-GRPO 是单条 trajectory 内的 step selection。前者解决 ``好轨迹别浪费''，后者解决 ``低熵步骤别乱训''。
  \item \textbf{和 DenseGRPO：} DenseGRPO 把 reward 变密，E-GRPO 把探索步骤变少但更有信息量。一个偏 reward attribution，一个偏 exploration attribution。
  \item \textbf{和 TempFlow-GRPO：} TempFlow-GRPO 也注意到 timestep 不应等权；E-GRPO 则从 SDE entropy 出发，给出更直接的 step merging 和 ODE/SDE hybrid sampling 机制。
\end{itemize}

\begin{intubox}
\textbf{如果组合使用：} 一个自然方向是 OP-GRPO 复用高 reward trajectories，但只复用 / 更新其中 E-GRPO 认为 high-entropy 或 consolidated informative 的 steps。这样可能同时减少采样成本和无效更新成本。
\end{intubox}

\subsection{10. 我认为它的局限}

E-GRPO 的思想很清楚，但也有几个边界：

\begin{enumerate}[leftmargin=1.5em]
  \item \textbf{Entropy proxy 不等于真实任务重要性。} 高熵 step 更可能产生差异，但差异大不一定等于对某个 reward 最关键；任务相关 credit 仍然需要 reward 侧方法补充。
  \item \textbf{阈值 \(\tau\) 仍是超参数。} 论文默认设为 2.2，并在附录做了阈值消融，但换模型、换 scheduler、换采样步数后可能需要重新校准。
  \item \textbf{主要验证在 FLUX.1-dev + HPD。} 它是否同样适用于视频、图像编辑、3D 或 unified multimodal generation，还需要更多证据。
  \item \textbf{reward hacking 没有被根治。} 论文自己也承认 reward design 仍是瓶颈；E-GRPO 能减少无效步骤优化，但不能保证 reward 本身完全符合人类偏好。
\end{enumerate}

\begin{figure}[htbp]
\centering
\includegraphics[width=0.95\linewidth]{egrpo/08_failure_cases.png}
\caption{E-GRPO 的 failure cases。这个图提醒我们：entropy-aware step selection 能提升 RL 更新效率，但不能替代更可靠的 reward model、prompt grounding 和视觉一致性约束。图源：\href{https://arxiv.org/html/2601.00423v1}{E-GRPO arXiv HTML Figure 8}。}
\label{fig:egrpo-failure-cases}
\end{figure}

\subsection{11. 读完后的记忆钩子}

\begin{summarybox}
\textbf{我最后会这样概括：} OP-GRPO 说 ``昂贵采样出来的好轨迹别只用一次''；E-GRPO 说 ``昂贵轨迹里的每一步也别都当成有效探索''。前者是 trajectory-level efficiency，后者是 timestep-level exploration efficiency。两篇连起来，正好构成 Flow-GRPO 采样效率支线的两个层级。
\end{summarybox}

\subsection{参考来源}

\begin{itemize}[leftmargin=1.5em]
  \item \textbf{论文主页：} \href{https://arxiv.org/abs/2601.00423}{E-GRPO: High Entropy Steps Drive Effective Reinforcement Learning for Flow Models}。
  \item \textbf{arXiv HTML：} \href{https://arxiv.org/html/2601.00423v1}{arXiv HTML version}。
  \item \textbf{论文配图：} \href{https://arxiv.org/html/2601.00423v1}{arXiv HTML figures}，已用于本文的 entropy analysis、pipeline、reward ambiguity、training curve、visual comparison 和 failure cases 图示。
  \item \textbf{代码仓库：} \href{https://github.com/shengjun-zhang/VisualGRPO}{shengjun-zhang/VisualGRPO}。
\end{itemize}
